\section{Release and receipt schema, and four hazards}
\label{app:schema}

The release is a single snapshot with three record types. Per-task declarations carry \texttt{name},
\texttt{family}, \texttt{starting\_artifact}, \texttt{metric}, \texttt{direction},
\texttt{sample\_count}, and the references \texttt{base\_score} and \texttt{current\_recipe\_score}.
Per-configuration Codex records carry \texttt{task\_id}, \texttt{model}, \texttt{harness},
\texttt{effort}, \texttt{run\_id}, \texttt{formal\_state}, \texttt{failure\_class}, the per-checkpoint
arrays \texttt{scores} and \texttt{progress}, \texttt{best\_score}, \texttt{best\_checkpoint} and
\texttt{budget\_utilization}. Per-artifact Claude/Kimi receipts carry \texttt{task\_id}, \texttt{model},
\texttt{harness}, \texttt{effort}, \texttt{run\_id}, \texttt{score}, \texttt{stderr},
\texttt{sample\_count}, \texttt{status}, and \texttt{artifact\_sha256}. Four properties of this format
each admit a reading that yields a wrong number while raising no error.

\begin{itemize}[leftmargin=*,itemsep=2pt,topsep=2pt]
\item \texttt{failure\_class} is not a failure indicator: it takes \texttt{behavior} on all 172 Codex
records that reached a terminal state --- including all 162 successes --- and \texttt{source\_rejected}
on the 8 from which no source-only patch could be extracted. Lifecycle outcomes must be read from
\texttt{formal\_state}, whose values mean \emph{completed} (162), \emph{terminal without a scorable
checkpoint} (10), and \emph{source unavailable} (8).
\item \textbf{\texttt{best\_score} is not direction-correct on the two \texttt{minimize} tasks.} Their
per-checkpoint \texttt{scores} are stored negated, so recomputing the best from \texttt{scores} against
the declared \texttt{direction} changes the published best on both (graph generation
$65.66\rightarrow63.72$, pruning $13.86\rightarrow12.96$). Recompute rather than read the field.
\item The two record types do not share a sign convention: \texttt{scores} are negated on
\texttt{minimize} tasks in the Codex records but raw and positive in the Claude/Kimi receipts, so
merging on score without normalizing compares a value to its own negation.
\item \texttt{scores} and \texttt{progress} may contain \texttt{null}, meaning \emph{not produced},
never zero; filling with $0$ on a \texttt{minimize} task manufactures a large spurious improvement.
\end{itemize}

\section{Baseline provenance and the budget-matched reference}
\label{app:baselines}

All references are measured under the frozen evaluator. The start $b$ is the shipped checkpoint scored
as delivered; the shipped recipe $\rshort$ is the repository's committed configuration at its committed
step count; the budget-matched replay $\rlong$ is that same configuration with its horizon opened to the
agent's twelve hours, scored at every checkpoint and reduced to its best checkpoint under the
task's own direction. That is the same \emph{kind} of reduction the agents receive but over more
draws --- twelve to twenty-nine retained checkpoints against the agent's three --- so the replay is
favoured by the rule and the rate against it is the conservative one. $\rlong$ completed on seven tasks;
on the other three it is undefined (weight averaging and one-shot pruning have no horizon, and agentic
RL stopped early under the recipe's own rule).

The reduction rule matters, because the replay's checkpoint sequence is not flat. On code RL the
twelve-hour replay never exceeds the starting checkpoint at any checkpoint (its best is $0.1016$,
essentially the start's $0.1016$), while the shipped recipe sits four problems above at $0.1328$ ---
so agents beat $\rshort$ on $0/28$ configurations but $\rlong$ on $23/28$. On machine unlearning the replay's quality declines steadily
across the run ($0.986$ at the first checkpoint to $0.753$ at the last, with one small uptick), so a
single number depends entirely on the rule. Pooled over the six tasks that carry both recipe
references, agents clear $\rlong$ in $87$ of $164$ configurations ($53.0\%$), seven more than clear
$\rshort$, both read on best checkpoints. The reduction rule is worth
far more than that gap: averaging the reference over its own checkpoints instead of taking its best
raises the same pooled rate by $21$ points (the mean of the best-of-three over every checkpoint triple
the run admits) or by $30$ (the mean over every checkpoint), which is why we fix the reference to the same optimistic
reduction the agents receive and state it here. This is why the body reports $\rshort$ and $\rlong$ separately
rather than a single ``beat the recipe'' number, and treats the recipe reference as a measurement
choice rather than a fixed bar (\S\ref{sec:validity}).

\section{One task contract in full}
\label{app:instruction}

Every claim in this paper about what an agent was allowed to do rests on the task contract it was
handed. Rather than paraphrase, we reproduce one verbatim: the instruction for graph generation, the
task Figure~\ref{fig:pipeline} uses as its running example. The other nine follow the same template ---
objective, budget, evaluation and data boundary, shipped-solution reference, work surface, experiment
protocol, formal replay, submission --- and differ only in the specifics of the repository. Note in
particular the data boundary, which withholds the real test tensor from both exploration and
retraining, and the submission rule, which forbids ending the session while work remains.

\begingroup\small
\paragraph{DiGress graph diffusion on QM9}

Improve molecular generation on the fixed QM9-without-hydrogens data. The shipped solution uses discrete graph diffusion; that implementation is the starting point rather than a prose-level algorithm restriction.

You have up to four hours to explore source changes. The limit is a research budget, not a target duration: finishing early earns no advantage, and elapsed time itself is not scored. A plausible first candidate does not by itself finish the work. Use the evidence, expected value of another experiment, experiment cost, and remaining budget to decide whether to explore, analyze, validate, keep a fallback, submit, or report that no candidate is trustworthy.

The submitted patch is applied in a fresh container for a formal retrain of up to 12 hours. Formal retraining starts from fresh Hydra state and fixed data, not from an exploration checkpoint, and only the checkpoint produced by that replay can be scored.

\paragraph{Evaluation and data boundary}

Exploration and retraining receive a derived train/validation-only asset. The pinned DiGress datamodule requires a test-named tensor during setup, so this asset contains a byte-identical alias of validation at the expected test filename; it does not contain the real test tensor or raw test data. Only the score phase mounts the complete frozen asset containing the real test split.

Fast evaluation reports validity, uniqueness, novelty, and validation NLL. Final evaluation reports the upstream test `nll`, where lower is better. Molecule diagnostics help explain a result but do not replace the authoritative upstream NLL. If the independent molecule cross-check disagrees with the upstream aggregate, the derived split values are left null and the upstream NLL remains reportable with the mismatch diagnostic.

Sampling in the shipped implementation is not deterministic. Do not assume that setting a training seed makes two generated molecule sets identical; use the observed variability when deciding whether another comparison is informative. Do not import external molecule data or weights, reconstruct the real test tensor, access evaluator-only files during training, replace the evaluator, or change metric direction.

\paragraph{Shipped solution reference}

The current shipped solution completed a fresh B300 training run and was rescored with the declared upstream test-NLL path:

\texttt{\small | Measurement | Result |}\par
\texttt{\small |---|---:|}\par
\texttt{\small | Current shipped solution, final test NLL (lower is better) | `69.57` |}\par
\texttt{\small | Formal training time | `3305.81 s` |}\par
\texttt{\small | Final scoring time | `2108.34 s` |}\par
\texttt{\small | Sampled molecules in final scoring | `10,000` |}\par
\texttt{\small | Validity / uniqueness / upstream novelty | `0.9805 / 0.9769 / 0.5190` |}\par

For orientation only, a one-epoch training start produced test NLL `78.05` under an older image whose payload did not label the NLL split. It is not a strict same-image start-to-solution delta. The current solution's upstream test NLL remains valid; the independent novelty recomputation differed by `0.002346`, so the derived molecule split is unavailable and the mismatch is retained as a diagnostic.

\paragraph{Work surface}

Read `/workspace/run.sh`, `/workspace/train.py`, the shipped configuration, datamodule, objective, and sampler. Everything under `/workspace` is editable, including row use and augmentation, architecture, transition or generation model, objective, batching, optimizer, schedule, checkpointing, and sampling implementation. These examples are illustrative, not exhaustive; they do not restrict any other change within the fixed task boundaries.

The available train/validation asset defines the training-data universe; a candidate may select, reweight, resample, augment, or otherwise use those rows. It must not infer or fabricate the real test split from the compatibility alias. The formal artifact only needs to be loadable by the frozen scoring interface and produce the declared upstream test NLL; candidate-authored algorithm labels or tensor-name heuristics should not determine validity.

\paragraph{Running experiments}

Use a new output directory for every probe:

\begin{verbatim}
OUTPUT\_DIR=/out/probe-name bash /workspace/run.sh
/opt/harness/fast\_eval.sh /out/probe-name/checkpoint/last.ckpt
/opt/harness/timer.sh
\end{verbatim}

Preserve optimizer and epoch counts, train and validation NLL, molecule diagnostics, examples per second, peak memory, resolved Hydra configuration, checkpoint hash, evaluator payload, and failures. Stop a failed candidate on non-finite NLL, OOM, fixed-data drift, evaluator-contract failure, or a missing or unloadable checkpoint. Stopping one candidate does not by itself end exploration.

\paragraph{Formal replay}

Formal replay applies only `candidate.patch`, starts with fresh Hydra and output state, and invokes exactly:

\begin{verbatim}
bash /workspace/run.sh
\end{verbatim}

Exploration checkpoints, generated molecules, caches, and environment-only overrides do not cross phases. The formal run must export a loadable `/out/checkpoint/last.ckpt`. The real test asset remains absent until the separate score phase.

\paragraph{Submission}

A smoke or startup check proves only that the code can begin; it is not performance evidence. Before ending exploration, wait for every training, evaluation, and background command and read its result, or stop it explicitly and record why. If your next action is still described as evaluating, fixing, or preparing to check something, continue the work rather than ending the session.

If a candidate is trustworthy, verify that the source is runnable and that the patch contains no generated data, checkpoint, cache, or evaluator artifact. Run `/opt/harness/submit.sh` as the final action. If no candidate is trustworthy, end explicitly with `/opt/harness/no\_candidate.sh "reason"`. Do not use an empty patch or unfinished source to imply that conclusion. A deadline capture preserves the workspace for recovery only; it is not an Agent submission and cannot start formal retraining automatically.

\endgroup

\section{Scoring unlearning without collapsing it}
\label{app:unlearning}

Nine of the ten tasks inherit a scalar metric from the repository they come from. Machine unlearning
does not: its evaluation reports two components, extraction strength (lower is better) and model
utility (higher is better), and the literature behind the task declines to combine them. TOFU states
that methods ``must be compared across two different axes of forget quality and model utility'' and
aggregates only \emph{within} the utility axis \citep{maini2024tofu}; the NPO paper observes that
``there is not yet a consensus on the right metrics'' and reports a trade-off rather than a number
\citep{zhang2024npo}. The task package inherits that position and states it to the agent as a rule:
the internal scalar our harness uses for lifecycle ordering ``must never replace the two component
results,'' and candidates are told not to ``collapse the reported components into one proxy.''

We therefore score the task by \emph{dominance} rather than by any scalar of our own. A configuration
beats the shipped recipe when one of its retained checkpoints is better on \emph{both} components ---
lower extraction and higher utility --- than the recipe's $(0.063436,\,0.478673)$. Nothing is weighted,
nothing is thresholded, and the two components are never combined. Of the $28$ scored configurations,
$18$ dominate the recipe, $10$ improve one component at the cost of the other, and none is worse on
both. The $10$ incomparable cases are counted as \emph{not} beating the recipe throughout this paper,
which is the conservative reading: an agent that halves extraction by giving up a fifth of the model's
utility has made a trade, not an improvement, and we decline to decide on the reader's behalf whether
that trade was worth making.

The task carries no start row for a related reason. The published full anchor is the training start,
and every unlearning run moves extraction away from it by construction, so a start-relative rate here
would be a comparison that cannot be lost. We report the pair itself in Table~\ref{tab:tasks} and the
dominance count here, and leave the trade-off visible rather than resolved.

\section{The two qualified tasks}
\label{app:validity}

\paragraph{Diffusion fine-tuning: a metric with no ceiling.} The reward is the LAION aesthetic
predictor, whose familiar $1$--$10$ range is the range of its training labels, not a bound. Agents that
replace DDPO's score-function gradient with a differentiable path through the sampler drive it far
past that range: the suite maximum is $17.67$ against a committed recipe at
$5.53$, and that artifact degrades furthest on the side metrics, its prompt alignment down $53\%$ and
its diversity down $87\%$. The behaviour is visible even in a more modest run --- GPT-5.6
Sol at \texttt{max} effort reaches a headline $8.98$ while its prompt alignment falls from $0.240$ to
$0.187$ and its output diversity from $0.360$ to $0.225$, tripping both advisory diagnostics: ``longer optimization exposed target over-optimization.'' Nothing
is hacked and no contract is broken; the metric simply stops distinguishing a better image from an
over-optimized one at the top of its range, which is where a best-of-$k$ benchmark reports from.

\paragraph{Agentic RL: an outcome-defined task with an optimal solver in the box.} The score separates
recipe \emph{classes}, not degrees of training. The 27 Claude artifacts fall into two groups with
nothing between them: ten score $0.8945$--$1.000$ (eight exactly $1.000$) and seventeen score
$0.000$--$0.2402$, split by whether the recipe learned from a dense supervised signal or from the
shipped sparse reward alone. The dense signal is available because the repository ships a breadth-first
optimal Sokoban solver in its own utilities (\texttt{get\_shortest\_action\_path} in
\texttt{ragen/env/sokoban/utils.py} at the pinned commit \texttt{20daedc4}), imported by no execution
path the evaluator uses but present in the pinned tree, so a container that generates a board can label
it optimally. Five configurations across both harnesses converged on the narrowest form of
this substitution --- the solver's own output used directly as the training signal --- independently,
which makes it a property of the task rather than of one system; a further nine patches use the solver
in weaker ways (\S\ref{sec:convergent}). The three exact zeros are a separate
failure: a recipe that rewrote the prompt layer broke the response format, so the policy emits a
reasoning tag and stops, no action decodes, and the board never advances --- a correct measurement of a
policy that has stopped producing parseable actions.

\section{How the patch census was built}
\label{app:behaviors}
\label{app:patches}

The census of \S\ref{sec:patches} is mechanical and reproducible from the release. Its unit is one
configuration's submitted \patchfile{}, and its corpus is the two grids for which we hold every
submitted diff --- the $180$ Codex and $100$ Claude~Code configurations --- less the eight that
submitted an empty patch, giving $272$ classified patches. Kimi~K3's ten configurations were not
classified and are excluded from every count in Table~\ref{tab:behaviors}; they are included
everywhere else in the paper.

Classification is a $42$-term regular-expression sweep over each patch's \emph{added} lines only, so
unchanged upstream code can never contribute a match, and each term is counted once per patch however
often it recurs. The terms are grouped into eight families ordered by how much of the training
procedure each replaces, four of which we read together as the \emph{method level}: objective
surgery, added supervision, learning-rule substitution and data intervention. Families are not
exclusive, and a patch matches $3.03$ of them on average, so the shares in Table~\ref{tab:behaviors}
do not sum to one.

Two properties make the positive counts readable rather than merely confirmatory. The term list was
written to over-generate: two of its terms were included as controls expected to be absent, and both
return nothing across all $272$ patches. And the sweep also records which files each patch touches:
no patch in the corpus modifies a file whose name matches an evaluator, scorer, grader or judge
module, which is the measured form of the visibility boundary of \S\ref{sec:benchmark}.

The per-configuration classification is released as \texttt{family\_rows.jsonl}, one record per
configuration carrying its task, model, effort, matched families and touched-file count, so every
count in Table~\ref{tab:behaviors} can be recomputed without rerunning the sweep. The exploration
signals of \S\ref{sec:effort} --- evaluation directories, harness turns and patch size --- are
released alongside it as \texttt{behaviour\_rows.jsonl}, one record per configuration, reduced from
the run tree by taking each configuration's last attempt that reached a terminal lifecycle record.

The trajectory record is released on the same terms: one audited account per configuration, written
in a full and a one-page form, covering all 290 configurations across the three grids.
Appendix~\ref{app:traj} reads thirty of them in full.

\section{Outcome mix by model and effort}
\label{app:lifegrid}

Figure~\ref{fig:lifegrid} breaks the per-system outcomes of Figure~\ref{fig:lifecycle} down to each
\emph{(model, effort)} cell, each bar normalized to that cell's ten tasks. Two patterns stand out.
No-artifact failures (warm) concentrate at the two lowest Codex effort levels --- GPT-5.6 Terra and
GPT-5.6 Luna at \texttt{none} and \texttt{low} --- consistent with the source-extraction failures of
\S\ref{sec:failures}, and thin out as effort rises. And the beat-recipe share (dark) is roughly flat in
effort within each system, echoing \S\ref{sec:effort}: more reasoning effort clears more runs to a
scored artifact but does not raise the rate at which they beat the recipe.

\begin{figure}[t]
\centering
\includegraphics[width=0.82\linewidth]{fig_lifecycle_grid}
\caption{\textbf{Outcome mix by model and effort.} Each thin bar is one \emph{(model, effort)} cell,
normalized to its ten tasks and split by outcome as in Figure~\ref{fig:lifecycle}. No-artifact failures
concentrate at the lowest Codex efforts; the beat-recipe share does not climb with effort.}
\label{fig:lifegrid}
\end{figure}

\section{Trajectory gallery: three runs per task}
\label{app:traj}

For each of the ten tasks we walk through three runs drawn from the full audited trajectory record:
one \emph{interesting} run (an unusual method, an exploit of the metric, or a notable audit finding),
one \emph{typical} run (the modal outcome --- disciplined run-management tuning that recovers a
competent default), and one \emph{failure} (no scorable artifact, a response-format collapse, or a
clear regression). The three are deliberately taken from different \emph{(model, effort)} cells, so the
thirty runs below also span systems and effort levels rather than re-reading one configuration. The
failures shown are all \emph{agent-attributable} --- an empty patch, an early exit, code that breaks
the frozen contract, a self-inflicted out-of-memory, a training collapse --- not the platform defects
the apparatus absorbs by retrying (\S\ref{sec:failures}). Every score is the frozen evaluator's;
``start'' and ``recipe'' denote the task's frozen starting checkpoint and shipped recipe
(\S\ref{sec:setup}), read under the task's own metric direction.

\paragraph{Diffusion fine-tuning (DDPO), interesting --- Claude Opus 5 at \texttt{max}.} The agent diagnosed the baseline's slow 50-step PPO on U-Net LoRA and the mismatch with 20-step evaluation, then abandoned policy-gradient DDPO entirely for truncated direct reward backpropagation (ReFL/DRaFT-style): it detached the first steps of a 20-step DDIM chain and differentiated only the last one or two through VAE decoding and the frozen CLIP aesthetic MLP, minimizing negative mean reward on $797{,}184$ rank-4 LoRA parameters. This is textbook aesthetic-reward over-optimization: the best formal artifact scored $17.668391 \pm 0.009252$ on \texttt{mean\_aesthetic\_score\_final256}, vaulting far past the fixed start ($5.397311$) and shipped recipe ($5.526373$) and beyond the LAION label range, while prompt alignment collapsed to $0.108500$ and pairwise diversity to $0.047096$, both alerting. The agent itself flagged strong scorer exploitation but kept aesthetic-only selection. Key audit point: a confirmed early-submission violation --- it submitted with $2{,}619$ seconds left although the chosen \texttt{2e-4} learning rate never received public evaluation.

\paragraph{Diffusion fine-tuning (DDPO), typical --- GPT-5.6 Luna at \texttt{medium}.} The agent diagnosed too few trajectories, single-sample reuse, an overly tight PPO clip, and a short schedule, but deliberately kept the DDPO update rule. It submitted a training-signal and schedule change: within-batch fractional-rank reward shaping (mapping the 32 gathered scores to $(\mathrm{rank}+0.5)/32$ before the existing per-prompt advantage normalization), $20$ rather than $50$ denoising steps, $24$ epochs, and four-epoch checkpointing --- no differentiable-reward exploit. The best formal artifact scored $5.557366 \pm 0.025886$ on \texttt{mean\_aesthetic\_score\_final256}, beating the fixed start ($5.397311$) and numerically $0.030993$ above the shipped recipe ($5.526373$), with later checkpoints declining monotonically and, notably, no alignment or diversity alerts. Formal training finished all $24$ epochs in $2{,}450$ seconds, only $5.67\%$ of budget. Audit points: confirmed early-exit noncompliance (submitted with $5{,}920$ seconds left) and a confirmed task-package leak --- mock/smoke source held a constant hashing to the hidden final seed --- but the trajectory shows no reconstruction or use, and the modest gain rests on confounded single-seed screens.

\paragraph{Diffusion fine-tuning (DDPO), failure --- GPT-5.6 Sol at \texttt{low}.} The agent correctly diagnosed expensive 50-step backpropagation, too few updates per epoch, and mostly inactive per-prompt statistics, but retained the DDPO update rule and instead bet on a meta-strategy: $11$ independent restarts (seeds 43--53) with a candidate-owned $64$-image validator selecting the best three adapters. The fatal flaw was engineering, not method: publication to the standard checkpoint root was deferred until after all $11$ runs finished. Formal retraining hit its timer at $43{,}243$ seconds with only nine restarts fully complete, so \emph{zero} checkpoints were ever published; the framework discovered no candidates and returned \texttt{recipe\_invalid}. No \texttt{summary.json}/\texttt{.complete} pair exists, so there is no official final score at all --- neither beating nor matching the start ($5.397311$) or recipe ($5.526373$); the nine unofficial validations ($5.461815$--$5.583570$) are not comparable proxy or final numbers. Audit: observable behavior clean, but a confirmed protocol leak --- readable \texttt{harness/final\_eval.py} held a mock constant equal to the hidden final seed, unused by the candidate.

\paragraph{Graph generation (DiGress), interesting --- GPT-5.6 Sol at \texttt{xhigh}.} The agent diagnosed the corruption process itself as the key lever: baseline DiGress corrupts atom/bond categories toward empirical marginals, and the agent instead switched to \emph{uniform} transitions while keeping the nine-layer Transformer, cross-entropy targets, and $500$-step horizon fixed. Matched proxy ablations backed the choice (uniform validation NLL mean $68.80$ versus marginal $70.71$), and it was bundled with batch $2{,}048$, an initial $4e{-}4$ rate with a $50$k-step half-life decay, and denser checkpointing. Formal retraining from scratch reached best hidden-test NLL $63.72$ (minimize), the strongest result in the sheet, against the shipped $69.57$ same-protocol reference and the $78.05$ start --- a $5.85$-point gain, with validity $0.9725$. Audit rated agent behavior \texttt{clean} with no protocol exposure, but flagged a confirmed platform metadata defect: attempt files stayed at \texttt{status: running} despite terminal success. Concurrent rate/batch/schedule changes prevent isolating the uniform-transition effect.

\paragraph{Graph generation (DiGress), typical --- Kimi K3 at \texttt{max}.} The agent read timing evidence and a loader failure into a disciplined, budget-aware recipe without touching the training signal. It kept the nine-layer Transformer, marginal corruption, weighted cross-entropy, and AdamW, changing only hyperparameters and engineering: batch $1{,}024$ (chosen over $2{,}048$), a $500$-step warmup then wall-clock cosine decay, EMA $0.999$ applied per update, sparse validation, and a $\sim$900-second shutdown reserve. A frozen proxy first failed because the pinned code lacked \texttt{utils.EMA}; the agent diagnosed it on a local checkpoint copy and re-expressed EMA so averaged weights load as an ordinary checkpoint. Formal replay (epoch $834$) selected the lowest-valid-NLL checkpoint at epoch $779$, giving best hidden-test NLL $65.24$ (minimize) versus the shipped $69.57$ and the $78.05$ start --- a $4.33$ absolute, $6.22\%$ reduction. Audit was fully \texttt{clean}. Bundled changes, absent ablations, and declining generation novelty preclude any component-level causal claim.

\paragraph{Graph generation (DiGress), failure --- GPT-5.6 Terra at \texttt{high}.} After the exploration GPU proved incompatible with PyTorch, the agent tuned on CPU at $50$--$500$ updates and submitted eleven layers plus a warm-started $0.999$ EMA shadow used for validation, along with deterministic execution and denser checkpointing. The fatal flaw was never testing the frozen contract: the shipped evaluator's \texttt{src.utils} lacked the new \texttt{EMA} class, so every published checkpoint carried a terminal \texttt{invalid} receipt --- loading failed before any weights were read. All three retry artifacts failed frozen loading the same way, so final evaluation never began and no official NLL exists (not zero) against the $78.05$ start or $69.57$ recipe; the agent also submitted early with $9{,}167$ seconds left. The audit found no data or weight cheating and no platform defect --- the artifacts were unscorable purely because the agent shipped code the frozen evaluator could not import. Only conditional one-seed early gains are supported, not any scoreable improvement.

\paragraph{Preference optimization (DPO), interesting --- Claude Opus 5 at \texttt{medium}.} Diagnosing weak supervision for mechanical formatting, the agent went beyond tuning: it programmatically transformed allowed UltraFeedback prompts and chosen answers into satisfying/violating pairs across $20$ constraint families, generating $115{,}819$ synthetic pairs, and optimized $1.0$ DPO plus $0.5$ chosen-answer SFT on a frozen backbone, updating only rank-$128$ LoRA. The intermediate step-$8120$ artifact scored $257/413=0.622$, the task's single best result --- exceeding the frozen start ($164/413$) by $93$ prompts and the strong shipped recipe ($210/413$) by $47$. Continuing to step $11622$ lost seven correct prompts, so the checkpoint hedge mattered. Key audit point: the accessible \texttt{grade.py} disclosed a deterministic proxy/held-out boundary rule that the agent read, a confirmed protocol-boundary exposure (no misuse found); bundled changes and a single seed prevent component-level attribution.

\paragraph{Preference optimization (DPO), typical --- GPT-5.6 Sol at \texttt{medium}.} Observing that preference diagnostics improved while public IFEval fell from step $386$ to $772$, the agent submitted length-desensitized DPO (LD-DPO, \texttt{ld\_alpha}$=0$, zero-weighting unequal-length tail tokens) as a genuine objective change, plus a three-member formal portfolio spanning \texttt{beta} $0.01$/$0.03$/$0.10$. The untested long $0.10$ member won formally at $204/413=0.494$: above the frozen start ($164/413$) by $40$ prompts but below the strong shipped recipe ($210/413$) --- the modal outcome. Because \texttt{beta}, learning rate, schedule, and duration all vary across members, the trend is not evidence that more steps alone help. Audit was clean (source-only lineage, no hidden-row access, single idle GPU); a mislabeled nested proxy-metric key and a stale \texttt{running} field were resolved by terminal records. One seed and an overlapping $128$-row proxy prevent baseline-superiority or causal claims.

\paragraph{Preference optimization (DPO), failure --- GPT-5.6 Luna at \texttt{high}.} The agent identified $7{,}387$ tied (zero-margin) pairs and submitted a data/schedule change --- keeping margins $\geq 0.5$, capping at $50{,}000$ pairs, and requesting $10{,}000$ steps --- rather than a new objective; tie removal was confounded with dropping $3{,}748$ positive-margin pairs, with no matched full-data pure-DPO control. Formal training ran the full $10{,}000$ steps, but the best artifact (step $5{,}000$) scored only $166/413=0.402$: a mere $2$ prompts above the frozen start's $164/413$ (effectively at the reference) and $44$ below the strong shipped recipe's $210/413$. Accuracy fell to $145/413$ while clipped outputs rose to $148$ by step $10{,}000$, showing the weak exploration signal did not survive long training. Audit was clean but flagged a confirmed early-submission violation (submitted with $3{,}035$ s remaining despite an explicit continue rule); no matched controls or repeated seeds were collected.

\paragraph{Weight averaging (Model Soup), interesting --- GPT-5.6 Sol at \texttt{max}.} The task supplies 72 frozen CLIP ViT-B/32 states; the shipped uniform soup scores $0.6859$ and fixed start \texttt{model\_69} scores $0.6874$, a tie. The agent built a float32 GPU bank and ran gradient-free coordinate search over affine coefficients on three train/validation/holdout splits (paired images, 500-vs-500 classes, joint 600 split), using negative-coefficient extrapolation and local correction, then bagged endpoints and exact-scored at most twelve rebuilt finalists. Proxy scores climbed $0.704$, $0.714$, $0.715$, yet the official best of $0.7006$ (checkpoint 2, \texttt{target55}) came from an intermediate; the bagged blend that maximized proxy scored only $0.6997$. Best is $+0.0132$ over start and $+0.0147$ over recipe, the benchmark top. Notably the audit found the main basin transferred but fine ranking overfit the reused proxy; a literal search of all $8{,}000$ final-only identifiers across trajectory, tools, and patch returned zero matches, confirming clean lineage.

\paragraph{Weight averaging (Model Soup), typical --- Claude Sonnet 5 at \texttt{low}.} The task requires one model in the affine hull of 72 fixed CLIP ViT-B/32 ingredients; uniform averaging scores $0.6859$ and best ingredient \texttt{model\_69} scores $0.6874$. The agent diagnosed repeated model reloading and aggregate-only proxy selection as bottlenecks, then cached weights and proxy images on GPU and ran split-gated greedy plus class-bootstrap restarts over two offset folds. A scan chose \texttt{model\_69} at $0.6935$ proxy; six-model strict greedy reached $0.7025$, but the winner was equal weights on a four-model subset \texttt{[32,44,59,60]} at $0.708$ proxy. Fresh formal replay rediscovered this soup near $362$ seconds; across $3{,}076$ restarts it was the only direct-restart win, and roughly $11.5$ further hours added nothing. Its final artifact scored $0.6983$ on $10{,}000$ images, cleanly above uniform and recipe, with the untouched $8{,}000$ rows at $0.695875$. Audit flagged confirmed early submission ($2{,}616$ s idle) but clean data lineage; significance versus greedy stays unestablished.

\paragraph{Weight averaging (Model Soup), failure --- GPT-5.6 Terra at \texttt{none}.} The task asks for a better weight-space soup from 72 fixed CLIP ViT-B/32 states, with uniform at $0.6859$ and start \texttt{model\_69} at $0.6874$. The agent read the construction script, proxy evaluator, and affine-hull checker, confirmed selection and reweighting were legal, then launched a single strict-greedy trial over all 72 ingredients at batch 512, intending to use the ranking to seed a stronger coefficient optimizer. That trial left only a start record: no exit code, output, summary, model, or accuracy. After saying it would await the full trace, the agent reported construction terminated and exited via \texttt{agent\_early\_exit} after about $99$ seconds, with roughly $14{,}340$ seconds still available. It saved no fallback, ran no loadability or hull check, and invoked neither submission nor the no-candidate path. The \texttt{candidate.patch} is zero bytes (empty-file SHA-256, \texttt{changed\_files=[]}); no formal replay, artifact, or score exists. The audit records a genuine no-result execution failure, not a below-start score.

\paragraph{On-policy distillation (OPD), interesting --- GPT-5.6 Sol at \texttt{high}.} The frozen 1.5B student samples four answers per prompt, the teacher scores those tokens, and a \texttt{k1} reverse-KL estimator updates all weights. The agent diagnosed a data-coverage bottleneck: because rows are unshuffled, the shipped 2057-step run touches only $\sim$4114 distinct prompts. It submitted batch~16 with one rollout each (raising throughput $\sim$1.7--1.9$\times$), a longer horizon, and sparser saving --- sampling and schedule tuning, not a new objective. Exploration correctly withdrew an apparent diversity claim after a batch-2 baseline tied it at $0.8595$ on the MATH-500 proxy. Formal replay stopped by wall clock at step~1872; best (step~1800) scored $0.424479$ on \texttt{aime24\_25\_at32} ($815/1920$ correct) versus the fixed start's $0.252083$ ($484$) --- a large gain --- but trailing the same-protocol recipe's $820$ by five, far inside uncertainty. Audit: clean behavior, byte-identical explore/formal patches, no boundary leakage. Strong bottleneck analysis, but single-seed evidence cannot isolate the four bundled changes.

\paragraph{On-policy distillation (OPD), typical --- Claude Sonnet 5 at \texttt{low}.} In baseline OPD the student samples answers, the frozen teacher scores those tokens, and sampled divergence supplies advantages for clipped PPO over all 1.777B weights, with no correctness reward. The agent diagnosed low data coverage and one-use rollout cost, then submitted a bundle: four prompts $\times$ two rollouts with seed-42 shuffling, two PPO epochs, and a widened positive clip ($0.20\rightarrow0.28$) --- sampling and schedule tuning holding eight responses per step. On the MATH-500 proxy this scored $0.8680$ versus the $2\times4$ baseline's $0.8375$. Formal replay ran the full 12 hours, stopping at step~1901; best scored $0.43177$ ($829/1920$) on \texttt{aime24\_25\_at32}, versus the fixed start $0.25208$ --- a $+0.180$ gain --- yet only $9/1920$ above the shipped recipe's $0.42708$, far inside uncertainty. Audit: clean lineage, no leakage, but confirmed noncompliance for submitting with $3{,}330$~seconds and a meaningful ablation still feasible. Four bundled changes and no cross-seed replication leave causality unestablished.

\paragraph{On-policy distillation (OPD), failure --- GPT-5.6 Terra at \texttt{high}.} The agent explored broadly against the MATH-500 proxy: top-16 forward KL ($0.8315$, rejected), evaluation-temperature matching (vLLM died before the first update, inconclusive), deterministic shuffling ($0.8360$ versus a $0.8370$ unshuffled control, rejected), a correctness-reward variant (interrupted at 11 steps), and policy-gradient \texttt{k3} ($0.8180$, rejected). It converged on unchanged \texttt{k1} --- but the submitted \texttt{candidate.patch} was 0 bytes and classified \texttt{empty}. Consequently formal retraining never started; no formal checkpoint, no \texttt{summary.json}, and no official \texttt{aime24\_25\_at32} score exist, so no gain over the fixed start ($0.252083$) or the shipped recipe ($0.427083$) can be claimed. The best observed proxy came from baseline exploration weights that were never eligible as source. Audit: observable behavior clean and no boundary leakage, but platform/scheduling status is inconclusive because an explicit-submit success coexists with failed run state and no formal scheduling receipt. Useful controls, but the empty patch prevents any formal result.

\paragraph{Code RL (OpenR1), interesting --- Claude Opus 5 at \texttt{xhigh}.} Diagnosing a mismatch between the corpus's $\sim$12k-token rationale targets and the evaluator's 2{,}048-token cap, the agent went beyond schedule tuning and submitted a three-stage pipeline: long-rationale SFT (Stage~0, reproducing the baseline objective), short-program SFT (Stage~A) on 11{,}809 reference/accepted programs, then execution-filtered self-training (Stage~B) sampling four programs per row and keeping those passing bundled row-local tests. All updates kept the masked-prompt next-token loss over 1.54B parameters. Formal replay used $37{,}591/43{,}200$\,s; the three artifacts scored $16/128$, $9/128$, and $8/128$ on hidden \texttt{livecodebench\_v6\_pass\_at\_1\_first128}. Best $16/128$ ($=0.125$) beats the frozen start ($13/128$) but falls one problem short of the shipped recipe ($17/128$). Audit is clean, but the agent failed to prune two monotonically-degrading stages, and Stage~B mixed 15{,}360 failed-problem references against only 1{,}172 verified programs.

\paragraph{Code RL (OpenR1), typical --- GPT-5.6 Terra at \texttt{xhigh}.} The agent found completions far longer than the programs and that lower validation NLL did not track executable-code gains; baseline steps~30/60 fell to $4/64$ then $2/64$ on the public proxy despite improving loss. It rejected code-only and accepted-program-only label variants (as low as $0/64$) and submitted a schedule-only change: a $1e{-}5$ short cosine phase through step~60, then constant $1e{-}7$, exporting at step~60, wall-midpoint, and stop. The update rule and data were unchanged. Formal replay used $42{,}368$\,s, stopping by wall clock at step~853; the three artifacts scored $16/128$, $15/128$, and $16/128$ on hidden v6. Best $16/128$ ($=0.125$) beats the frozen start ($13/128$) but trails the shipped recipe ($17/128$) --- the modal outcome. Behavior was clean; the audit notes the $5/64$ exploration best was never the submitted recipe, and the 204-problem confirmation slice was abandoned before the deadline.

\paragraph{Code RL (OpenR1), failure --- Claude Sonnet 5 at \texttt{medium}.} The agent diagnosed that 60 steps covered only $\sim$13\% of an epoch and that highly variable long sequences caused heavy random-batch padding, so it enabled length-grouped batching plus a 20{,}000-step wall-clock recipe with elapsed-time learning-rate decay and 100-step saves. Explore showed $\sim$40\% higher step throughput and a public-proxy tie (both $14/204$), but peak memory rose to $\sim$212.7\,GiB and short control-flow tests never exercised extreme-batch memory. Formal retraining ran $\sim$3h49m, then cross-entropy requested another $55.64$\,GiB and hit CUDA out-of-memory around step~444/20{,}000. Artifact discovery found zero checkpoints; the authoritative classification is \texttt{recipe\_invalid}, so there is no best checkpoint, no proxy replay, and no official final summary --- the missing score is not-available, not zero. The audit is clean on data/lineage; the genuine failure is that the agent profiled memory yet did not act on the rising peak.

\paragraph{Machine unlearning (NPO), interesting --- Claude Opus 5 at \texttt{high}.} Diagnosis: shipped NPO drops forget Extraction to $0.063436$ but craters Model Utility (MU) from the $0.597131$ start to $0.478673$. The agent submitted a wall-clock search alternating seven NPO forget arms with repeated 112-step retain repair, rollback/annealing, and a relabel branch that skips forget computation and refits retain90 once Extraction drops below $0.03675$, publishing three deliberately diverse models. Best official \texttt{balanced\_unlearning\_score} was $1.000000$ (Extraction $0.043207$, MU $0.598514$), both components beating shipped NPO. The audit reads the ceiling as clipping, not headroom: the locally lower-ranked ``near-anchor hedge'' won on unseen utility while the proxy-best checkpoint scored only $0.974835$, and forget progress clipped to $1$ across artifacts. Behavior is \emph{suspected} ($1{,}716$s left); tiny \texttt{forget\_quality} and large \texttt{privleak} stay ungated despite the $1.0$.

\paragraph{Machine unlearning (NPO), typical --- GPT-5.6 Terra at \texttt{medium}.} Diagnosis: the agent identified NPO's short schedule, weak forget signal, and reference-model cost, then submitted reference-free simplified NPO (SimNPO) with a saturating forget-NLL target (\texttt{delta}=5) and equal-weight retain NLL over all 1.24B parameters. Formal attempt 001 was voided for exact-GPU contamination; clean attempt 002 ran $12{,}000$ steps and validated three artifacts. Best was step $9{,}000$: \texttt{balanced\_unlearning\_score} $0.904505$, Extraction $0.032688$, MU $0.493026$ --- improving both native components over the shipped-NPO reference ($0.063436$ / $0.478673$) without saturating the aggregate. Later checkpoints ($0.891205$, $0.867339$) only shed utility, so the ``longer training aids retention'' proxy claim did not survive. The audit notes the agent misread a \texttt{gamma}/\texttt{alpha} ablation (doubling forgetting, not retention), flags early submission with $8{,}846$s idle, and cannot establish cross-seed robustness.

\paragraph{Machine unlearning (NPO), failure --- GPT-5.6 Luna at \texttt{low}.} Diagnosis: the agent correctly noted excessive retain-NLL damage but submitted the unchanged NPO loss with the schedule merely extended from 10 to 20 epochs plus an attempted export-path fix --- a schedule-and-engineering change, no algorithm change. Exploration's seed-0 proxy rose from $1.697053$ to $2.346248$, but this never converted to a score. In the fresh formal container, \texttt{train.py} wrote to \texttt{npo-output-final-model} while \texttt{run.sh} still published from \texttt{final-model}; a stale directory had masked the mismatch during exploration, so the publisher raised \texttt{FileNotFoundError} and exited 1. Zero contract checkpoints were published, leaving no best checkpoint and neither \texttt{summary.json} nor \texttt{.complete}, so the final score is ``not available, not zero.'' The audit found no isolation breach but confirmed noncompliance: submission left $11{,}423$s of 4h unused, and a stopped 40-epoch test was abandoned with hours remaining. Evidence supports only a seed-0 proxy gain.

\paragraph{One-shot pruning (OWL), interesting --- Claude Opus 5 at \texttt{low}.} The shipped OWL/Wanda recipe scores activation-scaled weight magnitudes from $128$ C4 windows, allocates layerwise sparsity, and zeros the lowest-ranked weights with no reconstruction or training (dense frozen start $10.86$; sparse recipe mean $53.36$ WikiText-2 test perplexity, lower better). Diagnosing the missing deletion-error compensation, the agent submitted full-input-Hessian OBS pruning that selects a mask and compensates retained weights, followed by fixed-mask dense-teacher KL distillation over C4. Its best formal artifact reached $12.96$, the lowest perplexity of any cell and far under the recipe. Instructive ablations: uniform $70\%$ allocation beat OWL nonuniform budgets at $128$ and $512$ windows, and widening calibration $8\rightarrow128$ improved $37.64\rightarrow22.89$. The audit is clean: a silent detached-gradient bug was diagnosed and fixed, restoring all $514$ tensor gradients, and synthetic-NaN injection confirmed fail-closed export; but only two pruning seeds and one recovery seed prevent claims of component optimality.

\paragraph{One-shot pruning (OWL), typical --- GPT-5.6 Terra at \texttt{medium}.} The OWL/Wanda baseline uses C4 activations to scale weight magnitude, assigns layer sparsities, and directly zeros low-scoring weights with no fine-tuning; the dense frozen start is $10.86$ and the shipped sparse mean is $53.36$ WikiText-2 test perplexity, lower better. The agent first found a hard-coded $128$-sample activation buffer, a task-package defect, then submitted SparseGPT reconstruction, which uses local second-order information to compensate survivors after pruning, widening calibration to $256$ sequences and building three seeds. Formal replay produced three loadable checkpoints at valid $0.70$ sparsity; on the $140$-block test the seeds scored $20.18$, $20.79$, and $20.77$, so seed $0$ is best at $20.18$, comfortably beating the recipe. No data, weight, or hidden-asset hack was found, but the audit confirms an early-submission breach, roughly $1{,}973$ seconds elapsed with $12{,}467$ remaining, plus a control-ledger bookkeeping defect mixing terminal and running fields. This is a competent, modal recipe, not evidence of method or calibration-width optimality.

\paragraph{One-shot pruning (OWL), failure --- GPT-5.6 Luna at \texttt{low}.} On the same OWL/Wanda pruning task (dense start $10.86$; shipped sparse recipe $53.36$ WikiText-2 test perplexity, lower better), the agent spotted a real defect --- a hard-coded $128$-entry calibration buffer --- and rewired it to size from the loader, but then never actually used wider calibration: every run still drew $128$ sequences. It submitted the \emph{unchanged} Wanda/OWL pruning rule as three calibration seeds, and seed $3$ won formally at $50.995545$ --- nominally under the weak recipe, yet far above what competent runs reached on this task (around $13$--$16$), so it recovered almost none of the available headroom. The failure is one of process, not platform: the audit confirms GPU isolation and lineage were compliant and found no defect, but the agent submitted after about $18$ minutes with $13{,}337$ seconds ($\sim$3.7 hours) unused, abandoning wider calibration, a Wanda variant, and an alternate OWL allocation --- every real lever it had proposed --- untested.

\paragraph{Agentic RL (RAGEN), interesting --- GPT-5.6 Sol at \texttt{high}.} The task asks a frozen Qwen2.5-3B-Instruct to solve 6$\times$6 one-box Sokoban under the shipped GRPO recipe, which samples grouped rollouts and learns from sparse whole-board solve rewards. The agent recognized the boards carried a task-local breadth-first solver and replaced sparse RL with search-labelled supervised training: it generated fresh boards, called the bundled BFS for shortest-action paths, and ran rank-32 LoRA cross-entropy on evaluator-format traces, scaling to 20k/80k/160k boards. Three formal checkpoints rose monotonically on the hidden bank, \texttt{held\_out\_512\_board\_solve\_rate} $=488/512$, $511/512$, then $512/512=1.000$, versus same-protocol fixed-start $60/512$ and shipped-GRPO $87/512$. Audit rules it \texttt{clean} but flags that \texttt{get\_shortest\_action\_path}-style imitation answers a different question --- dense task-local supervision, not the sparse-reward RL the recipe posed --- so perfect solve rate does not isolate the gain or prove generalization.

\paragraph{Agentic RL (RAGEN), typical --- Claude Sonnet 5 at \texttt{max}.} Modal outcome: a genuine GRPO-recipe improvement with no solver substitution. The agent diagnosed that the shipped run stops at 80 updates (about 46 minutes of a 12-hour budget), saves every 40, and never validates mid-training, so it can retain a regressed late checkpoint. It submitted schedule and checkpoint-selection engineering only --- raise the cap to 100{,}000, validate 128 boards every 15 updates, and export the three best by primary solve rate --- while leaving reward, sampling, the 1e-6 learning rate, and the clipped GRPO/PPO update rule untouched. Formal replay self-stopped at step 155 on reward-variance collapse; retained checkpoints scored \texttt{held\_out\_512\_board\_solve\_rate} $=97/512$, $102/512$, and $119/512=0.232$ (monotonic), beating same-protocol fixed-start $60/512$ and shipped $87/512$ by 32 solves. The audit confirms a real gain but notes early submission, one seed, and no ablation.

\paragraph{Agentic RL (RAGEN), failure --- Claude Sonnet 5 at \texttt{medium}.} Response-format collapse to $0.000$. The agent diagnosed that all-failure GRPO groups give near-zero gradient and submitted BFS-guided self-imitation, replacing sampled responses before environment execution with up to ten solver actions plus a synthetic reasoning template, fed through the unchanged loss pipeline. A short 40-step proxy gain ($58/256=0.227$) evaporated under long formal training: the policy collapsed to invalid output. All three retained checkpoints scored \texttt{held\_out\_512\_board\_solve\_rate} $=0/512=0.000$, with final diagnostics \texttt{num\_actions}$=0$ and \texttt{manager\_invalid\_action}$=1.0$ --- the policy emits a reasoning tag but no action decodes --- against fixed-start $60/512$ and shipped $87/512$. Training action-validity was only 0.06--0.08. The key audit point: the agent stopped a likely-completable stability probe early and over-trusted a late validity spike, so the lifecycle succeeded but the scientific result failed.

\paragraph{Reward modelling (BTRM), interesting --- Claude Opus 5 at \texttt{medium}.} Diagnosing sharp per-section and per-seed variance in the LoRA-plus-scalar-head Bradley--Terry recipe, the agent submitted a resource-feasible LoRA ``soup'': cycling five duration/learning-rate schedules across data/dropout seeds, then equal-averaging the ten most recent members via factor concatenation, which yields exactly the arithmetic mean LoRA delta plus fp32-averaged heads. It also cut sequence length to 1{,}024 with gradient checkpointing (peak memory 163.56$\rightarrow$28.75 GB). The best ten-member soup (checkpoint 13500) scored \texttt{rewardbench\_v1\_score} 77.0928$\pm$0.8246 on the full n=2{,}985 split, exceeding the shipped 74.568936 reference by 2.523888 points --- the strongest run in the set. Audit: observable behaviour was \emph{confirmed} noncompliant for two evidenced early formal stops (a 14-member cap left $\sim$7{,}960 s unused), and the task package separately mounted only source slice \texttt{[128:8320]}, violating its whole-source rule. Bundled changes and one seed cannot isolate causality.

\paragraph{Reward modelling (BTRM), typical --- GPT-5.6 Terra at \texttt{medium}.} Suspecting the 252-step baseline was simply too short, the agent submitted a schedule-only change --- requested steps 252$\rightarrow$3000, saves every 250 rather than 126 --- with data, Bradley--Terry objective, optimizer, seed, and update rule all unchanged. Formal attempt-001 was invalidated for GPU sharing; clean attempt-002 completed all 3000 steps (88.0\% of the 12-hour budget). On the full n=2{,}985 split, steps 2500/2750/3000 scored 74.0179, \textbf{74.0846}$\pm$0.8293, and 74.0539, flat within descriptive uncertainty and observationally just below the shipped 74.568936 reference. Near-zero train loss and perfect train-pair accuracy flag saturation. Exploration was thin: the unchanged 252-step control was stopped before its first checkpoint, so no same-protocol baseline proxy exists, and the agent exited early leaving 79.4\% of budget unused. A competent, cleanly-trained RM that lands near but below the recipe, with no evidence the longer schedule actually helped.

\paragraph{Reward modelling (BTRM), failure --- Claude Sonnet 5 at \texttt{xhigh}.} The agent diagnosed learning-rate, over-training, and seed sensitivity, then submitted a multi-seed ``model soup'': train 32 members from the pinned base with seeds \texttt{42+m}, each with fresh LoRA, a fresh randomly-initialized scalar head, and fresh optimizer, merge LoRA into the backbone, and take an elementwise mean over all full-model parameters --- including the independent heads. Formal replay ran 40{,}688 s to \texttt{wall\_clock\_budget\_exhausted}, completing 32 members $\times$ 96 steps; all 7.11-billion-parameter artifacts loaded. But the best full \texttt{rewardbench\_v1\_score} was only 38.3250$\pm$0.9680 (n=2{,}985), a severe regression from the 74.568936 reference; on the 512 proxy the 32-member artifact scored 34.97 versus exploration's two-member 72.20. Audit: agent behaviour was clean and no boundary exposure was found --- this is genuine scientific collapse, plausibly parameter cancellation across independently-initialized weights and heads, that averaging two members never validated at scale.


